
\subsubsection{Spurious Correlation and Group Robustness}

Spurious correlations arise when features correlated with labels in training data do not hold in deployment. Sagawa et al. (2019) formalized this as a group robustness problem, introducing Group DRO to optimize worst-group accuracy (WGA). Group DRO achieves approximately 91\% WGA on Waterbirds but requires group annotations during training.

Subsequent work removes this requirement. JTT \citep{liu2021jtt} trains an initial ERM model, identifies misclassified samples as likely minorities, and upweights them in a second training stage, closing approximately 75\% of the gap to Group DRO on Waterbirds (achieving ~86\% WGA). DFR \citep{izmailov2022feature} demonstrates that ERM already learns good features---the problem is the classifier head---and achieves 97\% WGA by retraining only the last layer on a balanced validation set.

\subsubsection{Early Training Dynamics}

SPARE \citep{Yang2023} exploits simplicity bias for early spurious correlation identification, achieving +21.1\% WGA improvement while being $12\times$ faster than prior methods. LA-SSL \citep{zhu2023lassl} observes that minority samples learn slower than majority samples and uses inverse learning speed for sampling weights.

This work differs from these approaches by analyzing per-sample loss trajectories as continuous discriminative signals, rather than using fixed thresholds (SPARE) or aggregate learning speed (LA-SSL).

\subsubsection{Positioning}

\begin{table}[htbp]
\centering
\resizebox{\columnwidth}{!}{%
\begin{tabular}{llll}
\toprule
Method & Signal Type & Group Labels & Training Stages \\
\midrule
Group DRO & Worst-group loss & Required & 1 \\
JTT & Binary misclassification & Validation only & 2 \\
DFR & Feature reweighting & Balanced val set & 1.5 \\
SPARE & Fixed epoch threshold & None & 1 \\
This work & Loss magnitude at early epoch & None & 1 \\
\bottomrule
\end{tabular}
}
\end{table}

